\begin{tikzpicture}
\begin{axis}[
    ybar,
    title={公司 2023—2024 各季度净利润},
    ylabel={净利润（万元）},
    xlabel={季度},
    symbolic x coords={2023Q1, 2023Q2, 2023Q3, 2023Q4, 2024Q1, 2024Q2, 2024Q3, 2024Q4},
    xtick=data,
    x tick label style={font=\scriptsize, rotate=45, anchor=east},
    ymin=-100, ymax=350,
    enlarge x limits=0.15,
    every node near coord/.append style={font=\scriptsize, fill=none, draw=none, inner sep=1pt},
    legend style={at={(1.03,0.5)}, anchor=west, draw=black, fill=white},
    bar width=15pt
]
% 正负值区分色：两个 addplot 的 x 集合必须完全相同（缺失值补 0），
% 再配合 bar shift=0pt 让两组柱子回到同一个分类位置而不并排。
% 补 0 的柱子高度为 0、视觉上不存在，并用显式空标签 [] 抑制其数值标注。
%
% 反面写法（不可取）：把正值放一个 addplot、负值放另一个，两组 x 互不重叠。
% 那样 pgfplots 会把它们当分组柱的两个系列 → xtick=data 只取到一个系列的分类
% （另一半 X 轴标签消失），且两个系列各横向偏移半个柱宽、柱位与分类错开。
\addplot[fill=green!60!black, bar shift=0pt, nodes near coords, point meta=explicit symbolic]
    coordinates {
    (2023Q1,120)[120] (2023Q2,0)[] (2023Q3,210)[210] (2023Q4,0)[]
    (2024Q1,260)[260] (2024Q2,180)[180] (2024Q3,0)[] (2024Q4,320)[320]
};
\addplot[fill=red!80, bar shift=0pt, nodes near coords, point meta=explicit symbolic]
    coordinates {
    (2023Q1,0)[] (2023Q2,-85)[-85] (2023Q3,0)[] (2023Q4,-40)[-40]
    (2024Q1,0)[] (2024Q2,0)[] (2024Q3,-30)[-30] (2024Q4,0)[]
};
\legend{盈利, 亏损}
\end{axis}
\node[anchor=north west, font=\scriptsize] at (current bounding box.south west) {数据来源：公司内部财务数据};
\end{tikzpicture}
